\section{Closed parameter bands and an original-current Radon calculus}
\label{sec:r12-closed-parameter-bands}

Fix the original finite extension $K/\mathbf Q_p$ and its original
absolute value. Let $U$ be an ordinary open in a good, Hausdorff,
boundaryless, quasi-smooth, pure $n$-dimensional $K$-analytic space.
All currents and tests in this section are original CLD objects.
In particular, an equality is required on every original compact
smooth complementary test, allowing every actual moment refinement.
No equality on an enlarged GK delta-test sheaf is asserted.

Let $f_1,\ldots,f_c$ be original analytic units, set
\[
 s_i=\log|f_i|,\qquad s=(s_1,\ldots,s_c):U\longrightarrow J\subset\mathbf R^c,
 \qquad
 \Omega_s=\bigwedge_{i=1}^c(d's_i\wedge d''s_i),
 \tag{PB1}
\]
where $J$ is a real open containing the image under consideration.
No rank, independence, submersion, saturation, or product hypothesis
on $s$ is imposed in this section. The original smooth functions
$s_i$ satisfy $Ds_i=0$ and their prime and double-prime differentials
are separately closed.

\begin{definition}[Parameter-annihilated currents]
\label{def:r12-parameter-annihilated}
A coefficient-Radon current $R$ on $U$ is \emph{$s$-annihilated} if
\[
                 d's_i\wedge R=d''s_i\wedge R=0
                         \qquad(1\le i\le c).
 \tag{PB2}
\]
An allowed $s$-background is an $s$-annihilated, original symmetric
middle-positive, separately $d'$- and $d''$-closed diagonal current.
The adjective describes a proved extra structure of a derived
current; it does not restrict the background $S$ in OPEN-GQ.
\end{definition}

For every arbitrary original allowed $(p,p)$-background $S$ with
$p+c\le n$, the derived current $\Omega_s\wedge S$ is an allowed
$s$-background. Indeed $\Omega_s$ is smooth, strongly positive and
separately closed, so its wedge preserves the original middle cone;
every conormal in (PB2) is repeated in this product. This uses no
regularization or replacement of $S$. It also covers all cases where
$\Omega_s$ happens to vanish because of moment-rank degeneracy.

\subsection{The Borel parameter algebra preserves closedness}

\begin{lemma}[Bounded smooth approximation for finitely many parameter measures]
\label{lem:r12-parameter-L1-approximation}
Let $\nu_1,\ldots,\nu_N$ be finite positive Radon measures on a
compact subset of $J$, and let $b:J\to\mathbf R$ be Borel and
bounded on a neighborhood of that subset. There are smooth real
functions $b_m$ on $J$, with one common bound, such that
\[
                   \sum_{j=1}^N\int|b_m-b|\,d\nu_j\longrightarrow0.
 \tag{PB3}
\]
If $b\ge0$, the $b_m$ can be chosen nonnegative. No measure in this
statement is assumed to be nonatomic.
\end{lemma}

\noindent\textbf{Status: \texttt{EXTERNAL\_KNOWN}} as elementary
Radon regularity; the short proof fixes precisely the version used.

\begin{proof}
Replace the finite family by its sum. Truncate $b$ to its fixed
bound on a compact neighborhood and approximate it in $L^1$ by
bounded simple functions. Inner and outer regularity approximate
each indicator by a continuous function between zero and one.
Continuous compactly supported functions on $J$ can be uniformly
approximated by smooth functions by convolution after a compact
cutoff. Truncation with a smooth approximation to the scalar
clipping map keeps a common bound and, when required, nonnegativity.
A diagonal choice proves (PB3).
\end{proof}

\begin{theorem}[Closed Borel bands and differential commutation]
\label{thm:r12-closed-parameter-bands}
Let $R$ be an allowed $s$-background. For every locally bounded
Borel $b:J\to\mathbf R$, the intrinsic coefficient multiplication
$b(s)R$ is separately closed and $s$-annihilated. It is middle-positive
if $b\ge0$. In particular, for every Borel $E\subset J$,
\[
                            R_E:=\mathbf1_{s^{-1}(E)}R
 \tag{PB4}
\]
is a genuine positive separately closed original current.

More generally, let $u$ be an original locally bounded Borel scalar
such that $T=D(uR)$ is coefficient-Radon. Then $T$ is $s$-annihilated
and, on all original tests,
\[
                         D\bigl(u\,b(s)R\bigr)=b(s)D(uR).
 \tag{PB5}
\]
Thus (PB5) applies in particular to every frozen locally bounded
classical BT insertion into $R$.
\end{theorem}

\noindent\textbf{Status: \texttt{NEW\_PROVED}} for the original
closed-band calculus. The coefficient-Radon nature and positivity of
the frozen bounded BT insertion are \texttt{INTERNAL\_INPUT}.

\begin{proof}
First take $b$ smooth. Its differentials are combinations of the
$d's_i$ and $d''s_i$. Consequently
\[
 d'b(s)\wedge R=d''b(s)\wedge R=Db(s)\wedge R=0.
 \tag{PB6}
\]
The distributional product rule and separate closedness of $R$
give separate closedness of $b(s)R$. Conormality follows immediately
from (PB2). For a compact original test $\eta$ complementary to $T$,
the smooth identity
\[
 D(b(s)\eta)=b(s)D\eta+d'b(s)\wedge d''\eta
                    -d''b(s)\wedge d'\eta+Db(s)\wedge\eta
\]
and (PB6), still valid after scalar multiplication by $u$, imply
\[
 \langle b(s)D(uR),\eta\rangle
   =\langle uR,D(b(s)\eta)\rangle
   =\langle uR,b(s)D\eta\rangle.
 \tag{PB7}
\]
This is (PB5) for smooth $b$.

For Borel $b$, fix one arbitrary compact original test stage. Choose
finite coefficient-Radon frames controlling $R$ on the supports of
the test and its required derivatives. Push the total variations of
these finitely many coefficient measures forward by $s$; they are
finite Radon measures on the real parameter space. Apply
Lemma~\ref{lem:r12-parameter-L1-approximation}. The resulting
$b_m(s)R$ converge to $b(s)R$ against every pairing needed for the
closedness equation of that test. Distributional differentiation is
continuous, so closedness passes to the limit. Multiplication of the
already zero coefficient measures in (PB2) proves conormality.
Positive scalar multiplication proves middle positivity when $b\ge0$.

For (PB5), include in the same finite controlling family the
coefficient measures of $T$ on $\operatorname{supp}\eta$ and those
of $R$ on $\operatorname{supp}D\eta$. The common compact bound on
$u$ then gives convergence of both sides of (PB7). This proves
(PB5) without inferring convergence against a discontinuous weight
from weak current convergence alone. The approximating functions
may depend on the test: since that test was arbitrary, the resulting
intrinsic equality holds on the whole original test sheaf. No
countability assumption on the collection of all refinements is
needed.

Finally, separate closedness of $d's_i,d''s_i$ gives directly, as
distributions,
\[
 d's_i\wedge D(uR)=D(u\,d's_i\wedge R)=0,
 \qquad
 d''s_i\wedge D(uR)=D(u\,d''s_i\wedge R)=0.
 \tag{PB8}
\]
This proves the asserted conormality of $T$. All arguments involve
intrinsic original forms and coefficient measures, and hence are
unchanged after any actual refinement containing the chosen data.
\end{proof}

\begin{corollary}[A closed current-valued Radon resolution]
\label{cor:r12-current-valued-parameter-resolution}
The assignment $E\mapsto R_E$ is countably additive on Borel sets
in every original compact-test pairing. Its local coefficient
variation is controlled by that of $R$. Every $R_E$ is separately
closed and positive. For all permitted bounded classical tuples,
\[
 T_{R_E,k}(u_1,\ldots,u_k)
       =\mathbf1_{s^{-1}(E)}T_{R,k}(u_1,\ldots,u_k).
 \tag{PB9}
\]
If $E_1,\ldots,E_N$ form a finite Borel partition on the parameter
range under consideration, then $R=\sum_iR_{E_i}$ and every bounded
mixed BT product has the corresponding exact finite decomposition.
The same conclusions hold for countable partitions whenever a
compact test is fixed and the sums are interpreted in coefficient
variation.
\end{corollary}

\noindent\textbf{Status: \texttt{DIRECT\_COROLLARY}}.

\begin{proof}
Coefficient-Radon multiplication gives countable additivity and the
variation bound. The theorem gives positivity and closedness of each
band. Apply (PB5) repeatedly: every positive mixed tail is separately
closed, coefficient-Radon, and $s$-annihilated by (PB8). This proves
(PB9). Sum the resulting coefficient measures over a partition.
\end{proof}

\paragraph{What is and is not a fiber construction.}
The resolution (PB4) retains every Borel set, including individual
parameters and atoms. It does not discard an exceptional parameter
set. It is a current-valued measure on the parameter space, not a
construction of BT products over a completed residue field. Its band
currents live on the original $U$, are carried measure-theoretically
by $s^{-1}(E)$, and use every original test there. Their topological
supports are only asserted to lie in the closure of that carrier.
No conditional fiber current on a new analytic field is asserted.

The conormal hypothesis is essential. A general Borel restriction
of a closed current need not be closed. Moreover, the frozen R10
annular example shows that arbitrary Borel scalar coefficients
cannot commute with curvature insertion: for the unit Gauss point
$\xi$, $\mathbf1_{\{\xi\}}[U]=0$ but
$\mathbf1_{\{\xi\}}D((\log|T|)_+[U])=\delta_\xi$.
That example does not contradict (PB5), because $[U]$ is not
annihilated by $d'\log|T|$ and $d''\log|T|$.

\paragraph{Frozen inputs and exact scope.}
Only the frozen original coefficient-Radon calculus, smooth form
Leibniz rules, positivity under strongly positive smooth wedges,
and bounded arbitrary-background BT theorem are used. The R10
source \path{signed_descent/pws_differential_curvature_module.tex}
is used for its explicit warning against unrestricted Borel
coefficient linearity; its finite PS module is not enlarged
without the separate proof above. No R/B/A2 proof is re-audited.
